% !TEX program = xelatex
% 編譯：XeLaTeX（Overleaf 請把 Compiler 設成 XeLaTeX）
% 圖片請放在 figures/ 資料夾；找不到圖時會自動顯示紅字佔位框，不會編譯失敗。
% 紅字【待補：…】是我沒有資料、需要你確認或補上的地方。
\documentclass[a4paper,11pt]{article}
\usepackage[top=1.4cm,bottom=1.4cm,left=1.8cm,right=1.8cm]{geometry}
\usepackage{fontspec}
\usepackage{xeCJK}
\usepackage{graphicx,array,colortbl,xcolor,pgfkeys}
\usepackage[hidelinks]{hyperref}

% ---------- 字型（與簡歷、自傳同一套：中文標楷體、英文 Times New Roman） ----------
\IfFontExistsTF{Times New Roman}
  {\setmainfont{Times New Roman}}
  {\setmainfont{Liberation Serif}}
\IfFontExistsTF{DFKai-SB}
  {\setCJKmainfont[AutoFakeBold=2.5]{DFKai-SB}}
  {\IfFontExistsTF{BiauKai}
    {\setCJKmainfont[AutoFakeBold=2.5]{BiauKai}}
    {\IfFontExistsTF{TW-Kai}
      {\setCJKmainfont[AutoFakeBold=2.5]{TW-Kai}}
      {\setCJKmainfont[AutoFakeBold=2.5]{AR PL UKai TW}}}}
\xeCJKsetup{CJKecglue={\hskip 0.2em plus 0.1em}}

\pagestyle{plain}
\setlength{\parindent}{0pt}
\renewcommand{\arraystretch}{1.25}

% ---------- 顏色（與簡歷、自傳同一套） ----------
\definecolor{navy}{RGB}{31,73,125}
\definecolor{hdrfill}{RGB}{221,228,238}
\arrayrulecolor{navy!55}
\hypersetup{colorlinks=true,urlcolor=navy,linkcolor=navy}
\newcommand{\hd}{hdrfill}
\newcommand{\todo}[1]{\textcolor{red}{【待補：#1】}}

% ---------- 圖片：找不到檔案就顯示佔位框 ----------
% 佔位框尺寸會依 #1 裡的 width/height 自動縮放（而非固定 2.7cm×3cm）。
% 固定尺寸在排 3 張並排縮圖（如 shotC）時會比欄位還寬，溢出的量大到讓
% colortbl 算錯同一張表格裡其他列（甚至更早的列）的 \cellcolor 色塊寬度，
% 在表格裡留下一道沒上色的縫隙——問題不在縮圖本身，而在佔位框。
% 找不到檔案時的說明文字改用 \nolinkurl（hyperref／url 套件），可在
% 「/」「.」等字元處自然換行，避免檔名在窄欄位裡擠出框外。
\pgfkeys{/safeimg/.is family, /safeimg,
  width/.store in=\simgW, height/.store in=\simgH,
  keepaspectratio/.style={}}
\newcommand{\safeimg}[2]{%
  \IfFileExists{#2}{\includegraphics[#1,keepaspectratio]{#2}}%
  {\begingroup
    \def\simgW{2.7cm}\def\simgH{3cm}%
    \pgfkeys{/safeimg,#1}%
    \fbox{\parbox[c][\simgH][c]{\dimexpr\simgW-2\fboxsep-2\fboxrule\relax}{%
      \centering\footnotesize\textcolor{red}{【圖片待補】}\\%
      \centering\fontsize{5.5}{6.8}\selectfont\nolinkurl{#2}}}%
  \endgroup}}

% ---------- 表格寬度（4 欄） ----------
\newlength{\cwT}\newlength{\cwA}\newlength{\cwB}\newlength{\cwC}\newlength{\cwD}
\newlength{\cwAB}\newlength{\cwBCD}\newlength{\cwCD}\newlength{\cwALL}
\newcommand{\setwidths}{%
  \setlength{\cwT}{\dimexpr\linewidth-8\tabcolsep-5\arrayrulewidth\relax}%
  \setlength{\cwA}{0.20\cwT}\setlength{\cwB}{0.24\cwT}%
  \setlength{\cwC}{0.32\cwT}\setlength{\cwD}{0.24\cwT}%
  \setlength{\cwAB}{\dimexpr\cwA+\cwB+2\tabcolsep+\arrayrulewidth\relax}%
  \setlength{\cwBCD}{\dimexpr\cwB+\cwC+\cwD+4\tabcolsep+2\arrayrulewidth\relax}%
  \setlength{\cwCD}{\dimexpr\cwC+\cwD+2\tabcolsep+\arrayrulewidth\relax}%
  \setlength{\cwALL}{\dimexpr\cwA+\cwB+\cwC+\cwD+6\tabcolsep+3\arrayrulewidth\relax}}

% ---------- 專題欄位（key-value） ----------
\pgfkeys{/proj/.is family, /proj,
  heading/.store in=\pjHeading, name/.store in=\pjName,
  advisor/.store in=\pjAdvisor, students/.store in=\pjStudents,
  contest/.store in=\pjContest, host/.store in=\pjHost, award/.store in=\pjAward,
  contestlabel/.store in=\pjContestLabel, hostlabel/.store in=\pjHostLabel,
  awardlabel/.store in=\pjAwardLabel,
  desc/.code={\long\gdef\pjDesc{#1}},
  arch/.store in=\pjArch, archlabel/.store in=\pjArchLabel,
  wide/.store in=\pjWide, widelabel/.store in=\pjWideLabel,
  shotA/.store in=\pjShotA, shotB/.store in=\pjShotB, shotC/.store in=\pjShotC,
  shotlabel/.store in=\pjShotLabel, photofull/.store in=\pjPhotoFull, photoh/.store in=\pjPhotoH,
  photo/.store in=\pjPhoto, photolabel/.store in=\pjPhotoLabel,
  archh/.store in=\pjArchH, shoth/.store in=\pjShotH, wideh/.store in=\pjWideH}

\newcounter{pj}
\newcounter{pjcert}
\newcommand{\lab}[2]{\multicolumn{#1}{|>{\centering\arraybackslash}p{#2}|}}
\newcommand{\hdcell}[3]{\lab{#1}{#2}{\cellcolor{hdrfill}\color{navy}\bfseries #3}}
% 標籤文字都很短、不需要垂直置中，直接沿用已驗證沒問題的 \hdcell（p{} 欄位）。
\newcommand{\vhd}[3]{\hdcell{#1}{#2}{#3}}
\newcommand{\vcell}[3]{\multicolumn{#1}{>{\centering\arraybackslash}m{#2}|}{#3}}
% 儲存格內的列點：各行靠左、整塊置中，行間用 \\ 分隔
\newcommand{\blist}[1]{{\footnotesize\begin{tabular}[c]{@{}l@{}}#1\end{tabular}}}
\newcommand{\bi}{\textbullet\,}

\newcommand{\proj}[1]{%
  \def\pjHeading{}\def\pjName{}\def\pjAdvisor{}\def\pjStudents{}%
  \def\pjContest{}\def\pjHost{}\def\pjAward{}\gdef\pjDesc{}%
  \def\pjContestLabel{競賽名稱}\def\pjHostLabel{舉辦單位}\def\pjAwardLabel{獲獎榮譽}%
  \def\pjArch{}\def\pjArchLabel{架構圖}\def\pjWide{}\def\pjWideLabel{系統畫面}%
  \def\pjShotA{}\def\pjShotB{}\def\pjShotC{}\def\pjShotLabel{系統畫面}%
  \def\pjPhotoFull{}\def\pjPhotoH{}%
  \def\pjPhoto{}\def\pjPhotoLabel{活動照片／成果照片}\def\pjArchH{4.3cm}\def\pjShotH{4.2cm}\def\pjWideH{}%
  \pgfkeys{/proj,#1}%
  \ifnum\value{pj}>0 \clearpage\fi\stepcounter{pj}\setwidths
  \ifx\pjAdvisor\empty\def\pjAdvisorRow{}\else
    \def\pjAdvisorRow{\vhd{1}{\cwA}{指導教授}&\vcell{3}{\cwBCD}{\pjAdvisor}\\\hline}\fi
  \ifx\pjArch\empty\def\pjArchRows{}\else
    \def\pjArchRows{\hdcell{4}{\cwALL}{\pjArchLabel}\\\hline
      \multicolumn{4}{|c|}{\rule{0pt}{\dimexpr\pjArchH+2mm}\safeimg{width=0.95\cwALL,height=\pjArchH}{\pjArch}}\\\hline}\fi
  \ifx\pjWide\empty\def\pjWideRows{}\else
    \def\pjWideRows{\hdcell{4}{\cwALL}{\pjWideLabel}\\\hline
      \multicolumn{4}{|c|}{\rule{0pt}{\dimexpr\pjWideH+2mm}\safeimg{width=0.95\cwALL,height=\pjWideH}{\pjWide}}\\\hline}\fi
  \ifx\pjPhotoH\empty\let\pjPhotoH\pjShotH\fi \ifx\pjWideH\empty\let\pjWideH\pjArchH\fi
  \ifx\pjShotC\empty\def\pjShotW{0.46}\else\def\pjShotW{0.30}\fi
  \def\pjShotsFull{\rule{0pt}{\dimexpr\pjShotH+2mm}\safeimg{height=\pjShotH,width=\pjShotW\cwALL}{\pjShotA}%
    \ifx\pjShotB\empty\else\hspace{6pt}\safeimg{height=\pjShotH,width=\pjShotW\cwALL}{\pjShotB}\fi
    \ifx\pjShotC\empty\else\hspace{6pt}\safeimg{height=\pjShotH,width=\pjShotW\cwALL}{\pjShotC}\fi}%
  \def\pjShotRows{}%
  \ifx\pjShotA\empty
    \ifx\pjPhoto\empty\else
      \def\pjShotRows{\hdcell{4}{\cwALL}{\pjPhotoLabel}\\\hline
        \multicolumn{4}{|c|}{\rule{0pt}{\dimexpr\pjPhotoH+2mm}\safeimg{height=\pjPhotoH,width=0.9\cwALL}{\pjPhoto}}\\\hline}%
    \fi
  \else
    \ifx\pjPhoto\empty
      \def\pjShotRows{\hdcell{4}{\cwALL}{\pjShotLabel}\\\hline
        \multicolumn{4}{|c|}{\pjShotsFull}\\\hline}%
    \else
      \ifx\pjPhotoFull\empty
        \def\pjShotRows{\hdcell{2}{\cwAB}{\pjShotLabel}&\hdcell{2}{\cwCD}{\pjPhotoLabel}\\\hline
          \multicolumn{2}{|c|}{\rule{0pt}{\dimexpr\pjShotH+2mm}%
            \ifx\pjShotC\empty
              \safeimg{height=\pjShotH,width=0.46\cwAB}{\pjShotA}%
              \ifx\pjShotB\empty\else\hspace{4pt}\safeimg{height=\pjShotH,width=0.46\cwAB}{\pjShotB}\fi
            \else
              \safeimg{height=\pjShotH,width=0.30\cwAB}{\pjShotA}\hspace{3pt}%
              \safeimg{height=\pjShotH,width=0.30\cwAB}{\pjShotB}\hspace{3pt}%
              \safeimg{height=\pjShotH,width=0.30\cwAB}{\pjShotC}%
            \fi}&%
          \multicolumn{2}{c|}{\safeimg{height=\pjPhotoH,width=0.92\cwCD}{\pjPhoto}}\\\hline}%
      \else
        \def\pjShotRows{\hdcell{4}{\cwALL}{\pjShotLabel}\\\hline
          \multicolumn{4}{|c|}{\pjShotsFull}\\\hline
          \hdcell{4}{\cwALL}{\pjPhotoLabel}\\\hline
          \multicolumn{4}{|c|}{\rule{0pt}{\dimexpr\pjPhotoH+2mm}\safeimg{height=\pjPhotoH,width=0.9\cwALL}{\pjPhoto}}\\\hline}%
      \fi
    \fi
  \fi
  \par\vspace{0.5em}\noindent{\large\bfseries\color{navy} \thepj.\ \pjHeading}\par\vspace{0.3em}
  \noindent\begin{tabular}{|p{\cwA}|p{\cwB}|p{\cwC}|p{\cwD}|}\hline
    \vhd{1}{\cwA}{專題名稱}&\vcell{3}{\cwBCD}{\pjName}\\\hline
    \pjAdvisorRow
    \vhd{1}{\cwA}{參與學生}&\vcell{3}{\cwBCD}{\pjStudents}\\\hline
    \vhd{2}{\cwAB}{\pjContestLabel}&\vhd{1}{\cwC}{\pjHostLabel}&\vhd{1}{\cwD}{\pjAwardLabel}\\\hline
    \multicolumn{2}{|>{\centering\arraybackslash}m{\cwAB}|}{\pjContest}&\vcell{1}{\cwC}{\pjHost}&\vcell{1}{\cwD}{\pjAward}\\\hline
    \hdcell{4}{\cwALL}{簡介}\\\hline
    \multicolumn{4}{|p{\cwALL}|}{\setlength{\parindent}{2em}\pjDesc}\\\hline
    \pjArchRows
    \pjWideRows
    \pjShotRows
  \end{tabular}\par}

% ---------- 經歷（簡式，用於工讀／課程專題等篇幅較短的項目） ----------
\newcommand{\exper}[3]{\par\noindent{\bfseries\color{navy}#1}\hfill #2\par\vspace{2pt}#3\par\vspace{9pt}}

% ---------- 經歷卡片（key-value，樣式比照 \proj，用於實習／社團／服務等完整經歷） ----------
\pgfkeys{/exp/.is family, /exp,
  heading/.store in=\exHeading, org/.store in=\exOrg,
  role/.store in=\exRole, period/.store in=\exPeriod,
  super/.store in=\exSuper,
  desc/.code={\long\gdef\exDesc{#1}},
  arch/.store in=\exArch, archlabel/.store in=\exArchLabel,
  shotA/.store in=\exShotA, shotB/.store in=\exShotB,
  shotlabel/.store in=\exShotLabel,
  photo/.store in=\exPhoto, photolabel/.store in=\exPhotoLabel}

\newcounter{pjexp}

\newcommand{\expcard}[1]{%
  \def\exHeading{}\def\exOrg{}\def\exRole{}\def\exPeriod{}\def\exSuper{}%
  \gdef\exDesc{}%
  \def\exArch{}\def\exArchLabel{相關畫面／文件}%
  \def\exShotA{}\def\exShotB{}\def\exShotLabel{工作／活動畫面}%
  \def\exPhoto{}\def\exPhotoLabel{團隊照片／成果照片}%
  \pgfkeys{/exp,#1}%
  \ifnum\value{pjexp}>0 \clearpage\fi\stepcounter{pjexp}\setwidths
  \ifx\exArch\empty\def\exArchRows{}\else
    \def\exArchRows{\lab{4}{\cwALL}{\cellcolor{hdrfill}\color{navy}\bfseries \exArchLabel}\\\hline
      \multicolumn{4}{|c|}{\safeimg{width=0.68\cwALL,height=4.3cm}{\exArch}}\\\hline}\fi
  \def\exShotRows{}%
  \ifx\exShotA\empty
    \ifx\exPhoto\empty\else
      \def\exShotRows{\lab{4}{\cwALL}{\cellcolor{hdrfill}\color{navy}\bfseries \exPhotoLabel}\\\hline
        \multicolumn{4}{|c|}{\safeimg{height=3.6cm,width=0.9\cwALL}{\exPhoto}}\\\hline}%
    \fi
  \else
    \ifx\exPhoto\empty
      \def\exShotRows{\lab{4}{\cwALL}{\cellcolor{hdrfill}\color{navy}\bfseries \exShotLabel}\\\hline
        \multicolumn{4}{|c|}{\safeimg{height=3.6cm,width=0.46\cwALL}{\exShotA}%
          \ifx\exShotB\empty\else\hspace{6pt}\safeimg{height=3.6cm,width=0.46\cwALL}{\exShotB}\fi}\\\hline}%
    \else
      \def\exShotRows{\lab{2}{\cwAB}{\cellcolor{hdrfill}\color{navy}\bfseries \exShotLabel}&%
        \lab{2}{\cwCD}{\cellcolor{hdrfill}\color{navy}\bfseries \exPhotoLabel}\\\hline
        \multicolumn{2}{|c|}{\safeimg{height=3.6cm,width=0.46\cwAB}{\exShotA}%
          \ifx\exShotB\empty\else\hspace{4pt}\safeimg{height=3.6cm,width=0.46\cwAB}{\exShotB}\fi}&%
        \multicolumn{2}{c|}{\safeimg{height=3.6cm,width=0.92\cwCD}{\exPhoto}}\\\hline}%
    \fi
  \fi
  \ifx\exSuper\empty\def\exSuperRow{}\else
    \def\exSuperRow{\vhd{1}{\cwA}{指導／主管}&\vcell{3}{\cwBCD}{\exSuper}\\\hline}\fi
  \par\vspace{0.5em}\noindent{\large\bfseries\color{navy} \thepjexp.\ \exHeading}\par\vspace{0.3em}
  \noindent\begin{tabular}{|p{\cwA}|p{\cwB}|p{\cwC}|p{\cwD}|}\hline
    \vhd{1}{\cwA}{服務單位}&\vcell{3}{\cwBCD}{\exOrg}\\\hline
    \vhd{1}{\cwA}{職稱／角色}&\vcell{3}{\cwBCD}{\exRole}\\\hline
    \vhd{1}{\cwA}{服務期間}&\vcell{3}{\cwBCD}{\exPeriod}\\\hline
    \exSuperRow
    \lab{4}{\cwALL}{\cellcolor{hdrfill}\color{navy}\bfseries 內容}\\\hline
    \multicolumn{4}{|p{\cwALL}|}{\setlength{\parindent}{2em}\exDesc}\\\hline
    \exArchRows
    \exShotRows
  \end{tabular}\par}

% ---------- 附錄獎狀：一頁一張，置中加說明 ----------
\newcommand{\certpage}[2]{%
  \ifnum\value{pjcert}>0 \clearpage\fi\stepcounter{pjcert}%
  \begin{center}
    \safeimg{height=19cm,width=0.88\linewidth}{#1}\par\vspace{6pt}
    {\fontsize{10}{13}\selectfont #2}
  \end{center}}

% ---------- 節標題（navy 粗體 + navy 底線，取代預設 \section*） ----------
\newcommand{\gridsec}[1]{%
  \par\vspace{2pt}{\bfseries\fontsize{14}{17}\selectfont\color{navy}#1\par}%
  \vspace{1pt}\noindent\textcolor{navy}{\rule{\linewidth}{0.8pt}}\par\vspace{4pt}}

% =====================================================================
\begin{document}

\begin{center}
  {\LARGE\bfseries\color{navy} 專題與經歷佐證資料}\\[6pt]
  {\large 韓欣澄 Hsin-Cheng Han\quad|\quad 國立中央大學 資訊工程學系}\\[2pt]
  116 學年度碩士班甄試 \quad 其他有利審查資料
\end{center}
\vspace{2pt}\noindent\textcolor{navy}{\rule{\linewidth}{0.9pt}}

\gridsec{壹、競賽成果}

% ---------------------------------------------------------------- 1
\proj{
  heading={2025 客語 AI 應用黑客松創意發想大賽——冠軍},
  name={KaTch —— 客語影像辨識翻譯學習系統 APP},
  advisor={國立中央大學資訊工程學系 周立德 教授},
  students={陳品謙、張峪銓、韓欣澄},
  contest={2025 客語 AI 應用黑客松創意發想大賽},
  host={客家委員會},
  award={冠軍},
  desc={利用 Flutter 開發的跨平台 AI 學習 App，結合影像辨識與語言學習：使用者拍下身邊的物品，系統即時辨識並翻譯為客語。前端為 Flutter，後端以 Python Flask 提供 API，整合 Moondream 2 Vision-Language Model、Gemini 2.5 Flash 與客家委員會提供的客語語音與翻譯 API，支援四縣腔與海陸腔。功能包含客語辭典、拍照學單字、情境對話與單字記事本。本作品在 40 隊報名、10 隊進入決賽的競賽中獲得冠軍。},
  arch={figures/katch_arch.png}, archlabel={系統架構}, archh=4.8cm,
  shotA={figures/katch_ui1.png}, shotB={figures/katch_ui2.png}, shotC={figures/katch_ui3.png},
  shotlabel={系統畫面：單字詳情、圖片情境對話與客語對話}, shoth=4.2cm,
  photo={figures/katch_award.jpg}, photolabel={頒獎典禮（冠軍）}, photoh=4.2cm}

% ---------------------------------------------------------------- 4
\proj{
  heading={2025 土地銀行校園金融創意挑戰賽——金獎},
  name={綠易通（Green E Pass）—— 中小企業 ESG 評分與綠色融資輔助平台},
  students={陳品謙、劉庭妤、韓欣澄},
  contest={2025 土地銀行校園金融創意挑戰賽},
  host={臺灣土地銀行股份有限公司、財團法人台灣金融研訓院},
  award={金獎（全國第一）},
  desc={協助中小企業評估 ESG 表現，並對應到綠色融資的貸款利率級距。核心是透明的 rule-based 100 分評分架構；LLM（Groq LLaMA 3.3 70B 與 Gemini）只負責把改善建議寫成易讀的說明，不參與評分，讓結果可以被稽核。後端使用 Node.js 與 Express，提供 14 個 REST API，並可輸出內嵌中文字型的 PDF 報告。本作品自初賽報名超過 200 隊中脫穎而出，經初賽書審、複賽工作坊與決賽簡報後獲得金獎，為本屆競賽第一名。},
  arch={figures/greenepass_arch.png}, archlabel={ESG 評級與貸款利率級距對應}, archh=3.2cm,
  wide={figures/greenepass_ui2.png}, widelabel={評估結果與視覺化趨勢追蹤}, wideh=2.2cm,
  shotA={figures/greenepass_ui1.png}, shotlabel={OCR 智慧碳排查：上傳單據自動填入}, shoth=2.6cm}

% ---------------------------------------------------------------- 5
\proj{
  heading={2025 中央資工校內網路應用創意競賽——金牌獎},
  name={司馬古洛夫（SMARTGLOVE）—— 導航與感應智慧手套},
  advisor={國立中央大學資訊工程學系 周立德 教授（計算機網路課程）},
  students={陳品謙（組長）、陳聖勳、林昱呈、韓欣澄},
  contest={校內網路應用創意競賽},
  host={國立中央大學資訊工程學系},
  award={金牌獎},
  desc={現有導航多仰賴手機螢幕與語音提示，使用者行走或騎車時容易因低頭查看或配戴耳機而分心；需要刷卡進出門禁或搭電梯時，掏卡也常因雙手忙碌、下雨或戴手套而不便。SMARTGLOVE 將導航提示與感應功能整合進手套：手機端以 Flutter 開發，透過 Google Maps API 規劃路線並判斷轉彎方向，再經藍牙 BLE 將指令送到手套內的 ESP32 微控制器，以 LED 燈條即時顯示轉向提示。手套內並整合 MFRC522 RFID 模組，可作為門禁卡與電梯感應卡使用；另以手機加速度計與陀螺儀偵測角速度劇烈變化，判定跌倒時自動發送緊急求救訊息。期中原規劃以震動馬達提示轉向，實作後改為 LED 燈條呈現，並新增跌倒偵測功能。},
  arch={figures/netapp_arch.png}, archlabel={系統流程圖（App 導航 → BLE → ESP32 → LED）}, archh=4.0cm, shoth=2.8cm,
  shotA={figures/netapp_ui1.png}, shotB={figures/netapp_ui2.png}, shotlabel={實機 Demo：導航與緊急求救},
  photo={figures/netapp_photo.jpg}, photolabel={實機 Demo：RFID 門禁感應}}

% ---------------------------------------------------------------- 6
\proj{
  heading={AI CUP 基於時序資料之桌球戰術與結果預測競賽——佳作},
  name={桌球戰術預測：LightGBM + Markov ensemble},
  advisor={國立中央大學資訊工程學系 周立德 教授、蔡宗翰 教授},
  students={韓欣澄、陳俊瑋、陳聖勳},
  contest={AI CUP 基於時序資料之桌球戰術與結果預測競賽},
  host={教育部／台灣人工智慧學校},
  award={佳作},
  desc={預測桌球比賽中的戰術。以 LightGBM 搭配 Markov ensemble 建模，並使用 GroupKFold 驗證。過程中發現 public leaderboard 的分數因 data leakage（serverGetPoint 與舊測試資料重疊）而被灌水，因此改以 private-safe baseline 評估模型，避免被虛高的排行榜分數誤導。},
  arch={figures/tt_pipeline.png}, archlabel={模型流程與驗證設計}, archh=5.0cm}

% ---------------------------------------------------------------- 7
\proj{
  heading={HLB 智慧驅鳥系統——新工程教育計畫課程專題},
  name={HLB 智慧驅鳥系統（IoT + AI 風力發電機驅鳥）},
  advisor={國立中央大學 林子軒 教授},
  students={劉庭妤、郭沛婕、陳柏軒、韓欣澄},
  contestlabel={計畫／工坊}, contest={中央大學工學院苗圃工坊（114-1）},
  host={國立中央大學工學院},
  award={代表學校參展},
  awardlabel={成果},
  desc={起源於竹圍漁港的觀察：候鳥遷徙路線與離岸風機設置區域高度重疊，鳥類撞擊不僅危及自身，也會損傷風機葉片，影響發電效率與維運成本。系統採三層架構：感知層以 AI 視覺模組（HUB-8735\_ultra）偵測鳥類，並以 HC-SR04 超音波模組即時量測鳥類與風機的距離；核心控制器為 Raspberry Pi Pico 2W，依偵測結果驅動致動層的 WS2812B 智能燈條、蜂鳴器與風機馬達，並提供網頁伺服器供工作人員遠端監控。距離大於 65 公分為安全（綠燈）；65 公分以內進入警戒（黃燈並觸發蜂鳴器主動驅離）；小於 45 公分為危險，系統會自動強制風機停機，若強制停機未成功，工作人員可透過網頁介面手動緊急停機，並同步發送 LINE 警報通知。實測顯示 HC-SR04 在 0–50 公分範圍內平均誤差低於 5\%，系統自偵測到紅燈警示的總反應時間約 150 毫秒。本專題為教育部「新工程教育方法實驗與建構計畫」（C 類學院主題式課群建構計畫）之課程成果。在 115 年期中海報審查暨成果交流活動中，我以該計畫授予的「助教」身分代表團隊介紹本作品。},
  arch={figures/hlb_arch.png}, archlabel={系統架構}, archh=2.7cm, wideh=2.7cm, shoth=2.3cm,
  wide={figures/hlb_ui1.png}, widelabel={遠端監控網頁與 LINE 警報通知},
  shotA={figures/hlb_ui2.png}, shotlabel={AI 視覺辨識鳥類},
  photo={figures/hlb_expo.jpg}, photolabel={超音波測距硬體}}

% ---------------------------------------------------------------- 8
\proj{
  heading={2025 年生成式人工智慧創新應用國際競賽——佳作},
  name={Format-Fluid: Solving the Knowledge-Modality Gap（知識模態轉譯引擎）},
  advisor={國立中央大學資訊工程學系 周立德 教授},
  students={韓欣澄（國立中央大學資訊工程學系）、劉庭妤（國立中央大學土木工程學系）、韓松諺（國立臺北大學經濟學系）},
  contest={2025 International Generative Artificial Intelligence Innovation Application Competition（2025 年生成式人工智慧創新應用國際競賽）},
  host={長榮大學管理學院（主辦）；日本新潟青陵大學、越南阮必成大學、印尼愛蘭加大學（協辦）；中華民國教育部（指導）},
  award={佳作},
  desc={針對「知識模態落差」（Knowledge-Modality Gap）問題——知識多以固定的文字格式呈現，卻不符合每個學習者不同的認知習慣——設計 Format-Fluid 智慧自適應學習引擎，以多階段序列式 AI Chain 將複雜文本轉譯為不同認知模態。Stage 1 以 GPT-4 Turbo 擔任「知識萃取者」，將原文解構為結構化 JSON（實體、因果關係、關鍵論點），去除原文的語氣與修辭；Stage 2 依使用者選擇的模式分別以不同 AI agent 生成創意藍圖：Podcast 模式生成師生對話腳本（Alice 提問、Bob 講解），Comic 模式生成四格漫畫分鏡與圖像提示，Summary 模式生成階層式重點摘要與生活化類比；Stage 3 以 OpenAI TTS API（雙音色）與 DALL-E 3 API 將藍圖轉為最終的語音與圖像成果，前端即時呈現。},
  arch={figures/ff_podcast.jpg}, archlabel={系統畫面：Podcast 模式}, archh=4.4cm, shoth=3.2cm,
  shotA={figures/ff_comic.jpg}, shotB={figures/ff_summary.jpg}, shotlabel={系統畫面：Comic 模式與 Summary 模式}}

% =====================================================================
\clearpage
\setcounter{pj}{0}
\gridsec{貳、研究專題與論文}

% ---------------------------------------------------------------- 2
\proj{
  heading={研究專題・國科會大專生研究計畫},
  name={基於多模態圖譜推論的可溯源會議知識檢索系統},
  advisor={國立中央大學資訊工程學系 周立德 教授},
  students={韓欣澄、張峪銓},
  contestlabel={計畫名稱}, contest={國科會大專生研究計畫},
  hostlabel={補助單位}, host={國家科學及技術委員會},
  awardlabel={競賽成果}, award={\blist{\bi 資電學院專題競賽佳作\\ \bi InnoServe 參賽中\\ \bi 智策AI 參賽中}},
  desc={國際標準組織（如 W3C、IETF）累積了大量會議影片，沒有字幕時幾乎無法檢索。GRASP 以 OpenCV 做畫面遮罩與改良式 SSIM 擷取投影片，用 Phi-3.5-Vision 解析投影片內容、Whisper 產生帶時間戳記的逐字稿，再由 LLM 抽取 entity 與 relation，建立 Neo4j knowledge graph，並以 Louvain community detection 產生社群摘要，結合 LanceDB 向量檢索完成 GraphRAG 問答。每則回答都能追溯到影片秒數或投影片頁面。前端為 Chrome Extension（Manifest V3、Side Panel），後端部署於雲端，並以 Ragas 評估檢索與回答品質。系統已全部實作完成，目前同時報名 2026 InnoServe 資訊應用服務創新競賽與智策AI（資策會 AI 創新實戰大賽），兩項競賽均於參賽中。},
  arch={figures/grasp_arch.png}, archlabel={系統模組架構}, archh=1.8cm, wideh=1.8cm,
  wide={figures/grasp_graphrag.png}, widelabel={可溯源即時問答模組（GraphRAG 檢索與生成流程）},
  shotA={figures/grasp_ui1.png}, shotlabel={Chrome 擴充功能：即時問答、知識圖譜探索與逐字稿},
  photo={figures/grasp_ui2.png}, photolabel={完整使用頁面（線上會議同步側欄）}, shoth=2.0cm}

% ---------------------------------------------------------------- 3
\proj{
  heading={ESG 承諾驗證階層校準與解碼研究——NTCIR-19 論文投稿中},
  name={Hierarchical Projection as a Validity Layer for Multilingual ESG Promise Verification},
  advisor={國立中央大學資訊工程學系 周立德 教授、蔡宗翰 教授},
  students={韓欣澄、陳俊瑋、陳聖勳、黃士育},
  contestlabel={投稿會議}, contest={NTCIR-19 國際研討會 AI CUP Special Session},
  hostlabel={主辦單位}, host={日本國立情報學研究所（National Institute of Informatics, NII）},
  awardlabel={論文狀態}, award={審稿中},
  desc={任務是對 ESG 報告書段落同時預測四個有階層限制的輸出欄位（PS、VT、ES、EQ），合法組合共 17 種。我們以 shared-encoder four-head RoBERTa 作為固定的 base model，只改變 calibration 與 output rule，比較 independent argmax、deterministic projection、class-bias calibration 與 17-state valid-state decoding。資料為 2,000 個段落、來自 49 份報告；採 rotating three-way cross-fitting，並同時評估 document-disjoint 與 same-document 兩種切分，以四欄加權 macro-F1 為主要指標。研究發現 hierarchical projection 作為免重訓的 validity layer，可將中文資料的不合法組合比例由 12.55\% 降至零，whole-tuple accuracy 提升 0.035；並在英、法、日、韓四語系共 32 組實驗中，whole-tuple accuracy 全數提升。本研究是我們在 AI CUP 永續承諾驗證競賽獲得前標的方法之延伸，目前已整理為論文，投稿至 NTCIR-19 國際研討會 AI CUP Special Session，審稿中。},
  arch={figures/esg_hierarchy.png}, archlabel={任務階層：120 種組合中僅 17 種合法}, archh=4.2cm,
  wide={figures/esg_rules.png}, widelabel={七種決策規則（output rule × calibration）}, wideh=4.2cm}


% =====================================================================
\clearpage
\setcounter{pjexp}{0}
\gridsec{參、實習與實務專案}

\expcard{
  heading={Garmin（台灣）Web Application Services 團隊——Software Engineer Intern},
  org={Garmin（台灣）Web Application Services 團隊},
  role={Software Engineer Intern},
  period={2026/07 起（在職）},
  desc={參與亞太地區多套對外服務系統的開發與維運，工作橫跨 legacy 系統現代化、新功能開發、效能疑難排解與跨系統整合。在一套校園獎學金申請系統（新舊框架並存）上，從需求訪談、既有程式碼考古開始，獨立完成後台可編輯審核欄位、資格篩選邏輯、排程提醒信與跨系統批次整合等多個功能模組，並配合法遵需求調整申請表單欄位，完整走過程式碼審查、測試與上線前檢查流程。另外，針對一套 B2B 經銷商後台的批次資料上傳效能問題，運用 GitHub Copilot 輔助分析既有程式碼找出瓶頸並驗證改善成效；也參與公司級 Java 版本汰換計畫，將舊服務升級至最新 LTS 版本；並以 Vue 重構一套客服申請流程的前端頁面，修正過一個客服整合服務的商業邏輯問題，以及將一項內部 API 呼叫遷移至新版服務端點並協調跨環境的網路權限開通。完整經歷企業級版本控制、Code Review、測試報告撰寫到正式環境部署的流程。}}

\expcard{
  heading={國立中央大學 蔡宗翰教授實驗室——教育部 AI CUP 全國競賽平台開發與維運},
  org={國立中央大學 蔡宗翰教授實驗室},
  role={實習生},
  period={114 學年度（大三）起，迄今},
  super={蔡宗翰 教授},
  desc={負責\textbf{教育部 AI CUP 全國競賽平台}的開發與維運。平台以 Django、Flask、PostgreSQL 與 Docker 建置，支援數千名參賽者的報名服務，並與趨勢科技、玉山金控等企業協作。針對伺服器故障風險，自行建立自動監控腳本與 LINE Bot 故障通知，並利用實驗室閒置機器架設備援主機，搭配自動 DNS failover；維運期間平台系統異常率相較往年降低 70\% 以上。}}

\expcard{
  heading={帛瑀室內設計（Silky Jade Interior Design）——形象官網接案},
  org={帛瑀室內設計},
  role={獨立接案開發者},
  period={\todo{期間}},
  desc={為室內設計工作室獨立開發五頁式形象官網（首頁、關於、作品、服務、聯絡），負責需求溝通、版面設計到上線部署的完整流程。以 HTML、CSS、JavaScript 手刻前端，串接 EmailJS 處理聯絡表單寄送，並部署於 Vercel。網站已正式對外營運。\todo{補充：是否有 RWD／SEO／作品管理等值得一提的細節}},
  arch={figures/silkyjade_home.png}, archlabel={網站首頁},
  shotA={figures/silkyjade_1.png}, shotB={figures/silkyjade_2.png}, shotlabel={作品與聯絡頁面}}

% =====================================================================
\clearpage
\setcounter{pjexp}{0}
\gridsec{肆、領導、教學與服務}

\expcard{
  heading={國立中央大學資訊工程學系學會——學術部共同創辦人暨副部長},
  org={國立中央大學資訊工程學系學會},
  role={學術部共同創辦人暨副部長},
  period={114 學年度（大三）起，迄今},
  desc={大一起加入公關部，後續與同學共同建立學術部並擔任副部長。主導規劃 AWS Taipei、Trend Micro 等企業參訪，以及升學與留學經驗分享活動，累計超過 200 人次參與。}}

\expcard{
  heading={AWS Educate 雲端校園大使},
  org={AWS Educate},
  role={第九屆雲端校園大使},
  period={2026 年 10 月起（甫獲選上任）},
  desc={獲選為第九屆 AWS Educate 雲端校園大使，於校內推廣 AWS 雲端技術學習資源，服務期間尚在進行中。}}

\exper{中央大學環境工程研究所 網路與伺服器系統管理（工讀）}{大三（114 學年度）}
{負責研究所的網路與伺服器系統管理。}

\exper{資訊與社會服務 II 課程專題 組長}{2025 春季學期}
{帶領 7 人團隊製作 5 支資安知識影片，上架至中央大學資工系官方 YouTube 頻道。}

% =====================================================================
\clearpage
\setcounter{pjcert}{0}
\gridsec{伍、附錄：獲獎證明}

\certpage{figures/cert_katch.jpg}{2025 客語 AI 應用黑客松創意發想大賽——冠軍（客家委員會）}
\certpage{figures/cert_greenepass.jpg}{2025 土地銀行校園金融創意挑戰賽——金獎（臺灣土地銀行、台灣金融研訓院）}
\certpage{figures/cert_ncu_project.jpg}{\todo{其他獎狀，如校內網路應用創意競賽、生成式 AI 國際競賽等}}

\end{document}
